% =============================================================================
%  PROPOSTA DE FOMENTO À PESQUISA — CAPES / CNPq / FINEP / FAPESP
%  Template LaTeX com formatação ABNT para agências brasileiras de fomento
%  Compilar com: pdfLaTeX + biber (ou BibTeX) + pdfLaTeX × 2
% =============================================================================
%  ATENÇÃO: Substitua todos os campos entre [colchetes] pelos dados reais
%  de sua proposta antes de compilar.
% =============================================================================

\documentclass[
  12pt,
  a4paper,
  oneside,
  english,      % ordem: idioma principal (último) = português
  brazil
]{article}

% -----------------------------------------------------------------------------
% PACOTES ESSENCIAIS — FORMATAÇÃO ABNT
% -----------------------------------------------------------------------------
\usepackage[T1]{fontenc}
\usepackage[utf8]{inputenc}
\usepackage{lastpage}

% Margens ABNT (3 cm superior/esquerda, 2 cm inferior/direita)
\usepackage[
  top=3cm,
  bottom=2cm,
  left=3cm,
  right=2cm
]{geometry}

% Espaçamento 1,5 (ABNT)
\usepackage{setspace}
\onehalfspacing

% Citações e referências no estilo ABNT
\usepackage[alf]{abntex2cite}   % \cite{chave} → autor-data
\usepackage{abntex2abrev}       % abreviações ABNT

% Fonte Times New Roman ou similar (profissional)
\usepackage{times}

% -----------------------------------------------------------------------------
% PACOTES DE APOIO
% -----------------------------------------------------------------------------
\usepackage[brazil]{babel}      % hifenização e títulos em português
\usepackage{amsmath,amssymb}    % matemática
\usepackage{graphicx}           % figuras
\usepackage{color,xcolor}       % cores
\usepackage{booktabs}           % tabelas profissionais (\toprule, \midrule…)
\usepackage{array}              % colunas personalizadas em tabelas
\usepackage{tabularx}           % tabelas com largura ajustável
\usepackage{multirow}           % células que ocupam várias linhas
\usepackage{enumitem}           % listas personalizadas
\usepackage{hyperref}           % links internos e URLs
\usepackage{parskip}            % espaço entre parágrafos (sem indent)
\usepackage{float}              % posicionamento explícito [H]
\usepackage{textcomp}           % símbolos diversos
\usepackage{microtype}          % microtipografia

% Configuração do hyperref (links pretos, sem borda)
\hypersetup{
  colorlinks=true,
  linkcolor=black,
  citecolor=black,
  urlcolor=blue,
  breaklinks=true
}

% -----------------------------------------------------------------------------
% COMANDOS AUXILIARES
% -----------------------------------------------------------------------------
\newcommand{\lorem}{%
  Lorem ipsum dolor sit amet, consectetur adipiscing elit.
  Sed do eiusmod tempor incididunt ut labore et dolore magna aliqua.
  Ut enim ad minim veniam, quis nostrud exercitation ullamco laboris
  nisi ut aliquip ex ea commodo consequat.}

% -----------------------------------------------------------------------------
% INÍCIO DO DOCUMENTO
% -----------------------------------------------------------------------------
\begin{document}

% =============================================================================
% 1. CAPA / FOLHA DE ROSTO
% =============================================================================
\begin{titlepage}
  \begin{center}

    \vspace*{1.5cm}

    % Logo da instituição (opcional — comente se não houver)
    % \includegraphics[width=0.3\textwidth]{figuras/logo-instituicao.png}\\[1.5cm]

    {\large\textbf{[Nome da Instituição]}}\\[0.3cm]
    {\large\textbf{[Sigla da Instituição / Centro / Departamento]}}\\[2.0cm]

    {\Large\textbf{[Nome Completo do Coordenador(a)]}}\\[0.3cm]
    {\large[Titulação · Cargo · Programa de Pós-Graduação]}\\[2.0cm]

    % TÍTULO EM PORTUGUÊS
    \vbox{%
      \hrule height 1.5pt
      \vspace{0.4cm}
      {\huge\textbf{[Título do Projeto em Português]}}\\[0.3cm]
      {\large\textbf{[Título do Projeto em Português — subtítulo, se houver]}}
      \vspace{0.4cm}
      \hrule height 1.5pt
    }

    \vspace{1.0cm}

    % TÍTULO EM INGLÊS
    {\large\textit{[Project Title in English]}}\\[0.2cm]
    {\large\textit{[Subtitle in English, if applicable]}}

    \vfill

    % DADOS DA CHAMADA
    \begin{tabular}{rl}
      \textbf{Agência de Fomento:} & [CAPES / CNPq / FINEP / FAPESP] \\
      \textbf{Edital / Chamada:}   & [Número do Edital / Chamada] \\
      \textbf{Linha de Pesquisa:}  & [Nome da Linha] \\
      \textbf{Vigência:}           & [24] meses \\
      \textbf{Valor Solicitado:}   & R\$ [XXX.XXX,XX] \\
    \end{tabular}

    \vspace{1.0cm}

    % LOCAL E DATA
    {\large[Cidade] — [UF]}\\[0.2cm]
    {\large[Mês de 202X]}

  \end{center}
\end{titlepage}

% =============================================================================
% 2. IDENTIFICAÇÃO DO PROPONENTE E EQUIPE
% =============================================================================
\newpage
\section*{Identificação do Proponente e da Equipe}
\addcontentsline{toc}{section}{Identificação do Proponente e da Equipe}

\begin{tabular}{ll}
  \toprule
  \textbf{Coordenador(a):} & [Nome completo] \\
  \textbf{CPF:}            & [XXX.XXX.XXX-XX] \\
  \textbf{Instituição:}    & [Sigla] \\
  \textbf{Unidade:}        & [Centro / Instituto] \\
  \textbf{Departamento:}   & [Nome do departamento] \\
  \textbf{E-mail:}         & \href{mailto:[email@instituicao.br]}{[email@instituicao.br]} \\
  \textbf{Telefone:}       & [+55 (XX) XXXXX-XXXX] \\
  \textbf{ORCID:}          & \href{https://orcid.org/[XXXX-XXXX-XXXX-XXXX]}{[XXXX-XXXX-XXXX-XXXX]} \\
  \textbf{Currículo Lattes:} & \href{http://lattes.cnpq.br/[XXXXXXXX]}{[http://lattes.cnpq.br/XXXXXXXX]} \\
  \bottomrule
\end{tabular}

\vspace{0.5cm}

\begin{tabularx}{\textwidth}{lX}
  \toprule
  \textbf{Equipe executora:} & \\
  \midrule
  [Nome do Pesquisador 1] & [Instituição] — [papel / responsabilidade] \\
  [Nome do Pesquisador 2] & [Instituição] — [papel / responsabilidade] \\
  [Nome do Aluno de Doutorado 1] & [Programa de PG] — [bolsista / voluntário] \\
  [Nome do Aluno de Mestrado 1]  & [Programa de PG] — [bolsista / voluntário] \\
  [Nome do Aluno de IC 1]  & [Curso de graduação] — [bolsista / voluntário] \\
  % Acrescente linhas conforme necessário
  \bottomrule
\end{tabularx}

\vspace{0.3cm}

\begin{tabularx}{\textwidth}{lX}
  \toprule
  \textbf{Colaboradores externos:} & \\
  \midrule
  [Nome do Colaborador 1] & [Instituição estrangeira ou nacional] \\
  [Nome do Colaborador 2] & [Instituição estrangeira ou nacional] \\
  \bottomrule
\end{tabularx}

% =============================================================================
% 3. RESUMO (máx. 500 palavras)
% =============================================================================
\newpage
\section*{Resumo}
\addcontentsline{toc}{section}{Resumo}

% --- INÍCIO DO RESUMO (PORTUGUÊS) ---
\begin{otherlanguage}{brazil}
\begin{abstract}
  \noindent
  [Redija o resumo do projeto em português, com no máximo 500 palavras.
  O resumo deve contemplar: contextualização, problema de pesquisa,
  objetivos, metodologia, resultados esperados e impacto.
  Evite abreviaturas não consagradas e referências bibliográficas.]

  \vspace{0.8cm}

  \noindent\textbf{Palavras-chave:}
  [palavra-chave 1]. [palavra-chave 2]. [palavra-chave 3].
  [palavra-chave 4]. [palavra-chave 5].
\end{abstract}
\end{otherlanguage}

% --- ABSTRACT (INGLÊS) ---
\vspace{0.5cm}
\begin{otherlanguage}{english}
\begin{abstract}
  \noindent
  [Write the abstract in English, mirroring the Portuguese version.
  Maximum 500 words. Include: background, research problem, objectives,
  methodology, expected results, and impact.]

  \vspace{0.8cm}

  \noindent\textbf{Keywords:}
  [keyword 1]. [keyword 2]. [keyword 3].
  [keyword 4]. [keyword 5].
\end{abstract}
\end{otherlanguage}

% =============================================================================
% 4. INTRODUÇÃO E JUSTIFICATIVA
% =============================================================================
\newpage
\section{Introdução e Justificativa}

\subsection{Contextualização do Problema}

[Descreva o contexto geral da pesquisa. Apresente o estado da arte,
as lacunas identificadas na literatura e o problema concreto que a
proposta pretende investigar. Cite trabalhos relevantes (atente-se
ao formato autor-data das normas ABNT).]

De acordo com \cite{autor2023}, \lorem{}

\lorem{} \lorem{} Ademais, \citeonline{autor2020} demonstram que
\lorem{}

\subsection{Originalidade e Contribuição}

[Explique o caráter inovador ou original da proposta. Destaque:
\begin{itemize}
  \item O que diferencia este projeto de estudos anteriores;
  \item A abordagem teórico-metodológica original;
  \item A contribuição para o avanço do conhecimento na área;
  \item O potencial de gerar produtos, processos ou serviços inovadores.
\end{itemize}]

\subsection{Relevância Científica, Tecnológica e Social}

[Discuta a relevância da pesquisa sob três dimensões:
\begin{enumerate}[label=\textbf{\alph*)}]
  \item \textbf{Científica:} avanço teórico, preenchimento de lacunas,
        consolidação de linhas de pesquisa;
  \item \textbf{Tecnológica:} inovação, patentes, transferência de
        tecnologia, metodologias aplicáveis;
  \item \textbf{Social:} impacto na qualidade de vida, políticas
        públicas, formação de recursos humanos, inclusão produtiva.
\end{enumerate}]

\subsection{Alinhamento com a Chamada}

[Demonstre como o projeto se alinha aos objetivos, linhas temáticas
e prioridades do edital ao qual concorre.]

% =============================================================================
% 5. OBJETIVOS
% =============================================================================
\newpage
\section{Objetivos}

\subsection{Objetivo Geral}

[Apresente, em um único parágrafo, o objetivo central do projeto.
Inicie com um verbo no infinitivo: Analisar, Desenvolver, Investigar,
Propor, Avaliar, Compreender\ldots]

\begin{quote}
  \textit{[Exemplo: Analisar o impacto da implementação de políticas
  de economia circular na cadeia produtiva do setor de [setor] no
  Nordeste brasileiro, propondo um modelo de avaliação multicritério
  que integre indicadores ambientais, sociais e econômicos.]}
\end{quote}

\subsection{Objetivos Específicos}

Os objetivos específicos do projeto são:

\begin{enumerate}
  \item [Primeiro objetivo específico — verbo no infinitivo];
  \item [Segundo objetivo específico — verbo no infinitivo];
  \item [Terceiro objetivo específico — verbo no infinitivo];
  \item [Quarto objetivo específico — verbo no infinitivo];
  \item [Quinto objetivo específico — verbo no infinitivo];
  \item [Sexto objetivo específico — verbo no infinitivo].
\end{enumerate}

% =============================================================================
% 6. METODOLOGIA
% =============================================================================
\newpage
\section{Metodologia}

\subsection{Classificação da Pesquisa}

[Classifique a pesquisa quanto à:
\begin{itemize}
  \item Natureza: básica / aplicada;
  \item Abordagem: quantitativa / qualitativa / mista;
  \item Objetivos: exploratória / descritiva / explicativa;
  \item Procedimentos: experimental / bibliográfica / documental /
        estudo de caso / survey / quase-experimental / design science research.
\end{itemize}]

\subsection{Universo e Amostra}

[Descreva:
\begin{itemize}
  \item População-alvo ou universo da pesquisa;
  \item Critérios de inclusão e exclusão;
  \item Técnica de amostragem e tamanho da amostra;
  \item Cálculo do poder amostral, se aplicável.
\end{itemize}]

\subsection{Instrumentos e Procedimentos de Coleta de Dados}

[Detalhe os instrumentos e procedimentos:

\begin{description}
  \item[Revisão sistemática:] bases consultadas (Scopus, Web of Science,
        SciELO), string de busca, critérios de seleção, PRISMA.
  \item[Experimentos:] ambientes, equipamentos, softwares, protocolos.
  \item[Entrevistas / questionários:] roteiros, escalas validadas,
        pré-teste, plataforma de aplicação (Google Forms, SurveyMonkey).
  \item[Observação:] tipo (participante / não participante), registros.
  \item[Dados secundários:] fontes (IBGE, IPEA, DataSUS, bases
        patentárias), período de abrangência, variáveis selecionadas.
\end{description}]

\subsection{Procedimentos de Análise de Dados}

[Especifique as técnicas analíticas:

\begin{itemize}
  \item \textbf{Análise quantitativa:} estatística descritiva,
        testes de hipóteses (t de Student, ANOVA, qui-quadrado),
        correlação (Pearson/Spearman), regressão linear/múltipla/logística,
        análise multivariada, modelagem de equações estruturais,
        séries temporais, machine learning (especificar algoritmo);
  \item \textbf{Análise qualitativa:} análise de conteúdo (Bardin),
        análise temática, análise do discurso, grounded theory,
        softwares de apoio (NVivo, ATLAS.ti, MAXQDA);
  \item \textbf{Análise mista:} triangulação de métodos,
        validação cruzada.
\end{itemize}]

\subsection{Cronograma Metodológico}

[Relacione cada objetivo específico à(s) técnica(s) metodológica(s)
que serão empregadas para alcançá-lo. Vide Seção~\ref{sec:cronograma}
para o cronograma detalhado.]

\subsection{Aspectos Éticos}

\begin{itemize}
  \item \textbf{Pesquisa envolvendo seres humanos:} submissão ao Comitê
        de Ética em Pesquisa (CEP) via Plataforma Brasil; cumprimento
        da Resolução CNS nº 466/2012 e 510/2016;
  \item \textbf{Pesquisa envolvendo animais:} submissão à Comissão de
        Ética no Uso de Animais (CEUA);
  \item \textbf{Propriedade intelectual:} observância da Lei nº
        9.279/1996 (Propriedade Industrial) e Lei nº 10.973/2004
        (Inovação Tecnológica);
  \item \textbf{Confidencialidade:} garantia de anonimato dos
        participantes, consentimento livre e esclarecido (TCLE),
        guarda segura dos dados conforme LGPD (Lei nº 13.709/2018).
\end{itemize}

% =============================================================================
% 7. RESULTADOS ESPERADOS E IMPACTO
% =============================================================================
\newpage
\section{Resultados Esperados e Impacto}

\subsection{Produtos Esperados}

\begin{itemize}
  \item \textbf{Artigos científicos:} submissão de [N] artigos a
        periódicos indexados (Qualis A1 a B1) — estimativa de
        publicações por ano;
  \item \textbf{Dissertações e teses:} orientação de [N] trabalho(s)
        de conclusão de curso (TCC), [N] dissertação(ões) de mestrado
        e [N] tese(s) de doutorado;
  \item \textbf{Patentes / Propriedade intelectual:} depósito de
        pedido de patente de invenção, modelo de utilidade, registro
        de software ou desenho industrial (especificar);
  \item \textbf{Produtos técnico-tecnológicos:} softwares, bases de
        dados, metodologias, protocolos, relatórios técnicos;
  \item \textbf{Transferência de tecnologia:} acordos de cooperação,
        licenciamento, spin-off acadêmica;
  \item \textbf{Atividades de extensão:} cursos, workshops, eventos
        científicos, divulgação científica (mídias sociais, blog,
        entrevistas).
\end{itemize}

\subsection{Impacto Científico}

[Descreva como os resultados contribuirão para:
\begin{itemize}
  \item Avanço do conhecimento teórico na área;
  \item Formação de redes de pesquisa nacionais/internacionais;
  \item Fortalecimento dos programas de pós-graduação envolvidos;
  \item Aumento da produção científica qualificada.
\end{itemize}]

\subsection{Impacto Tecnológico e de Inovação}

[Descreva:
\begin{itemize}
  \item Potencial de inovação e patenteamento;
  \item Desenvolvimento de novas tecnologias ou processos;
  \item Adensamento tecnológico de cadeias produtivas.
\end{itemize}]

\subsection{Impacto Social e Econômico}

[Descreva os benefícios esperados para a sociedade:
\begin{itemize}
  \item Melhoria de indicadores sociais (saúde, educação, renda);
  \item Subsídios para políticas públicas;
  \item Geração de emprego e renda;
  \item Formação de recursos humanos qualificados.
\end{itemize}]

\subsection{Indicadores de Acompanhamento e Avaliação}

\begin{tabularx}{\textwidth}{lXX}
  \toprule
  \textbf{Indicador} & \textbf{Meta} & \textbf{Periodicidade} \\
  \midrule
  Artigos submetidos/publicados & [N] & Anual \\
  Orientações concluídas        & [N] & Anual \\
  Patentes depositadas          & [N] & Ao final \\
  Participações em eventos      & [N] & Anual \\
  Relatórios técnicos entregues & [N] & Semestral \\
  \bottomrule
\end{tabularx}

% =============================================================================
% 8. PLANO DE TRABALHO E CRONOGRAMA
% =============================================================================
\newpage
\section{Plano de Trabalho e Cronograma}
\label{sec:cronograma}

\subsection{Etapas do Projeto}

\begin{enumerate}
  \item \textbf{[Etapa 1 — Exemplo: Revisão sistemática da literatura]}
        \hfill \textit{[Meses 1–3]} \\
        Descrição detalhada da etapa. Insumos, atividades e entregáveis
        previstos.

  \item \textbf{[Etapa 2 — Exemplo: Coleta de dados em campo]}
        \hfill \textit{[Meses 4–10]} \\
        Descrição detalhada da etapa. Insumos, atividades e entregáveis
        previstos.

  \item \textbf{[Etapa 3 — Exemplo: Análise dos dados]}
        \hfill \textit{[Meses 11–16]} \\
        Descrição detalhada da etapa. Insumos, atividades e entregáveis
        previstos.

  \item \textbf{[Etapa 4 — Exemplo: Redação e submissão de artigos]}
        \hfill \textit{[Meses 17–22]} \\
        Descrição detalhada da etapa. Insumos, atividades e entregáveis
        previstos.

  \item \textbf{[Etapa 5 — Exemplo: Relatório final e prestação de contas]}
        \hfill \textit{[Meses 23–24]} \\
        Descrição detalhada da etapa.
\end{enumerate}

\subsection{Cronograma Físico}

\begin{table}[H]
  \centering
  \caption{Cronograma de execução do projeto (em meses).}
  \label{tab:cronograma}
  \footnotesize
  \begin{tabularx}{\textwidth}{p{4cm}*{12}{>{\centering\arraybackslash}p{0.9cm}}}
    \toprule
    \multirow{2}{*}{\textbf{Etapa / Atividade}} &
      \multicolumn{12}{c}{\textbf{Mês}} \\
      \cmidrule{2-13}
      & \textbf{1} & \textbf{2} & \textbf{3} & \textbf{4} & \textbf{5}
      & \textbf{6} & \textbf{7} & \textbf{8} & \textbf{9} & \textbf{10}
      & \textbf{11} & \textbf{12} \\
    \midrule
    Etapa 1 — Revisão         & \checkmark & \checkmark & \checkmark & & & & & & & & & \\
    Etapa 2 — Coleta          & & & & \checkmark & \checkmark & \checkmark & \checkmark & \checkmark & \checkmark & \checkmark & & \\
    Etapa 3 — Análise         & & & & & & & & & \checkmark & \checkmark & \checkmark & \checkmark \\
    % Adicione mais linhas para os meses 13-24 se a vigência for > 12 meses
    \bottomrule
  \end{tabularx}

  \vspace{0.3cm}
  {\small\textit{Nota:} O cronograma acima cobre os primeiros 12 meses.
  O detalhamento completo (24 meses) encontra-se em anexo.}
\end{table}

% Caso prefira um cronograma completo (24 meses):
\vspace{0.5cm}
\begin{table}[H]
  \centering
  \caption{Cronograma completo (24 meses).}
  \label{tab:cronograma-completo}
  \footnotesize
  \begin{tabularx}{\textwidth}{p{3.5cm}*{24}{>{\centering\arraybackslash}p{0.4cm}}}
    \toprule
    \textbf{Atividade} &
      \textbf{1} & \textbf{2} & \textbf{3} & \textbf{4} & \textbf{5}
      & \textbf{6} & \textbf{7} & \textbf{8} & \textbf{9} & \textbf{10}
      & \textbf{11} & \textbf{12}
      & \textbf{13} & \textbf{14} & \textbf{15} & \textbf{16} & \textbf{17}
      & \textbf{18} & \textbf{19} & \textbf{20} & \textbf{21} & \textbf{22}
      & \textbf{23} & \textbf{24} \\
    \midrule
    Revisão sist. & \checkmark & \checkmark & \checkmark & & & & & & & & & & & & & & & & & & & & & \\
    Coleta dados  & & & & \checkmark & \checkmark & \checkmark & \checkmark & \checkmark & \checkmark & \checkmark & & & & & & & & & & & & & & \\
    Análise       & & & & & & & & & \checkmark & \checkmark & \checkmark & \checkmark & \checkmark & \checkmark & \checkmark & & & & & & & & & \\
    Redação arts. & & & & & & & & & & & & & & & \checkmark & \checkmark & \checkmark & \checkmark & \checkmark & \checkmark & & & & \\
    Relat. final  & & & & & & & & & & & & & & & & & & & & & \checkmark & \checkmark & \checkmark & \checkmark \\
    \bottomrule
  \end{tabularx}
\end{table}

% =============================================================================
% 9. ORÇAMENTO
% =============================================================================
\newpage
\section{Orçamento}

\subsection{Resumo do Orçamento}

\begin{table}[H]
  \centering
  \caption{Resumo do orçamento solicitado por rubrica.}
  \label{tab:resumo-orcamento}
  \begin{tabularx}{\textwidth}{lXr}
    \toprule
    \textbf{Rubrica} & \textbf{Descrição} & \textbf{Valor (R\$)} \\
    \midrule
    Capital     & Equipamentos, material permanente         & [XX.XXX,XX] \\
    Custeio     & Material de consumo, serviços de terceiros & [XX.XXX,XX] \\
    Bolsas      & Bolsas de IC, Mestrado, Doutorado, Pós-Doc & [XX.XXX,XX] \\
    Diárias     & Viagens e diárias nacionais                 & [XX.XXX,XX] \\
    Taxas       & Taxas de publicação (APC), registro de patentes & [XX.XXX,XX] \\
    \midrule
    \textbf{Total} & & \textbf{R\$ [XXX.XXX,XX]} \\
    \bottomrule
  \end{tabularx}
\end{table}

\subsection{Detalhamento — Capital}

\begin{table}[H]
  \centering
  \caption{Detalhamento dos itens de capital.}
  \label{tab:orcamento-capital}
  \begin{tabularx}{\textwidth}{p{1.5cm}Xp{2cm}p{1.8cm}p{1.8cm}}
    \toprule
    \textbf{Item} & \textbf{Especificação} & \textbf{Justificativa} &
    \textbf{Quant.} & \textbf{Valor unit.} \\
    \midrule
    1 & [Descrição do equipamento 1] & [Necessário para…] & [N] & [XX.XXX,XX] \\
    2 & [Descrição do equipamento 2] & [Necessário para…] & [N] & [XX.XXX,XX] \\
    3 & [Descrição do equipamento 3] & [Necessário para…] & [N] & [XX.XXX,XX] \\
    \midrule
    \multicolumn{4}{r}{\textbf{Subtotal — Capital}} & \textbf{[XX.XXX,XX]} \\
    \bottomrule
  \end{tabularx}
\end{table}

\subsection{Detalhamento — Custeio}

\begin{table}[H]
  \centering
  \caption{Detalhamento dos itens de custeio.}
  \label{tab:orcamento-custeio}
  \begin{tabularx}{\textwidth}{p{1.5cm}Xp{2cm}p{1.8cm}p{1.8cm}}
    \toprule
    \textbf{Item} & \textbf{Especificação} & \textbf{Justificativa} &
    \textbf{Quant.} & \textbf{Valor unit.} \\
    \midrule
    1 & [Material de consumo 1] & [Necessário para…] & [N] & [XX.XXX,XX] \\
    2 & [Serviço terceirizado 1] & [Necessário para…] & [N] & [XX.XXX,XX] \\
    3 & [Passagens aéreas] & [Viagens para coleta/eventos] & [N] & [XX.XXX,XX] \\
    4 & [Diárias] & [N diárias × valor] & [N] & [XX.XXX,XX] \\
    \midrule
    \multicolumn{4}{r}{\textbf{Subtotal — Custeio}} & \textbf{[XX.XXX,XX]} \\
    \bottomrule
  \end{tabularx}
\end{table}

\subsection{Detalhamento — Bolsas}

\begin{table}[H]
  \centering
  \caption{Solicitação de bolsas.}
  \label{tab:orcamento-bolsas}
  \begin{tabularx}{\textwidth}{lXp{2cm}p{1.8cm}p{1.8cm}}
    \toprule
    \textbf{Nível} & \textbf{Perfil} & \textbf{Duração (meses)} &
    \textbf{Quant.} & \textbf{Valor total} \\
    \midrule
    Iniciação Científica & [Perfil do aluno] & [24] & [N] & [XX.XXX,XX] \\
    Mestrado   & [Perfil do aluno] & [24] & [N] & [XX.XXX,XX] \\
    Doutorado  & [Perfil do aluno] & [36] & [N] & [XX.XXX,XX] \\
    Pós-Doutorado & [Perfil do bolsista] & [12] & [N] & [XX.XXX,XX] \\
    \midrule
    \multicolumn{4}{r}{\textbf{Subtotal — Bolsas}} & \textbf{[XX.XXX,XX]} \\
    \bottomrule
  \end{tabularx}
\end{table}

\subsection{Plano de Aplicação}

\begin{table}[H]
  \centering
  \caption{Distribuição do orçamento por ano.}
  \label{tab:orcamento-anual}
  \begin{tabularx}{\textwidth}{lXXXX}
    \toprule
    \textbf{Rubrica} &
    \textbf{Ano 1 (R\$)} & \textbf{Ano 2 (R\$)} & \textbf{Total (R\$)} \\
    \midrule
    Capital     & [XX.XXX,XX] & [XX.XXX,XX] & [XX.XXX,XX] \\
    Custeio     & [XX.XXX,XX] & [XX.XXX,XX] & [XX.XXX,XX] \\
    Bolsas      & [XX.XXX,XX] & [XX.XXX,XX] & [XX.XXX,XX] \\
    Diárias     & [XX.XXX,XX] & [XX.XXX,XX] & [XX.XXX,XX] \\
    Taxas       & [XX.XXX,XX] & [XX.XXX,XX] & [XX.XXX,XX] \\
    \midrule
    \textbf{Total} & \textbf{[XX.XXX,XX]} & \textbf{[XX.XXX,XX]} &
    \textbf{[XX.XXX,XX]} \\
    \bottomrule
  \end{tabularx}
\end{table}

% =============================================================================
% 10. EQUIPE EXECUTORA
% =============================================================================
\newpage
\section{Equipe Executora}

\subsection{Coordenação}

\begin{tabularx}{\textwidth}{llX}
  \toprule
  \textbf{Nome} & \textbf{Titulação} & \textbf{Atribuição no Projeto} \\
  \midrule
  [Coordenador(a)] & [Dr.(a)] & Coordenação geral, orientação de
  teses/dissertações, redação de artigos, gestão de recursos. \\
  \bottomrule
\end{tabularx}

\subsection{Pesquisadores Colaboradores}

\begin{tabularx}{\textwidth}{llX}
  \toprule
  \textbf{Nome} & \textbf{Instituição} & \textbf{Atribuição no Projeto} \\
  \midrule
  [Colaborador 1] & [Instituição] & Co-orientação, análise de dados,
  experimentos. \\
  [Colaborador 2] & [Instituição] & Coleta de dados, curadoria,
  validação. \\
  \bottomrule
\end{tabularx}

\subsection{Alunos e Bolsistas}

\begin{tabularx}{\textwidth}{llX}
  \toprule
  \textbf{Nome} & \textbf{Nível} & \textbf{Atribuição no Projeto} \\
  \midrule
  [Aluno 1] & Doutorado & Tese vinculada, experimentos, análise. \\
  [Aluno 2] & Mestrado  & Dissertação vinculada, coleta de dados. \\
  [Aluno 3] & Iniciação Científica & Apoio técnico, tabulação. \\
  \bottomrule
\end{tabularx}

\subsection{Infraestrutura Disponibilizada}

[Liste a infraestrutura existente que será disponibilizada para o projeto:

\begin{itemize}
  \item \textbf{Laboratórios:} [nome do laboratório], [equipamentos principais];
  \item \textbf{Equipamentos multi-usuário:} [especificar];
  \item \textbf{Acesso a bases de dados:} Portal de Periódicos CAPES,
        Web of Science, Scopus, SciELO;
  \item \textbf{Oficinas / plantas piloto:} [descrever];
  \item \textbf{Apoio administrativo:} [secretaria, núcleo de projetos].
\end{itemize}]

\subsection{Contrapartida da Instituição Proponente}

[Descreva a contrapartida oferecida pela instituição, se exigida pelo
edital: espaço físico, recursos humanos, equipamentos já adquiridos,
isenção de taxas, etc.]

% =============================================================================
% 11. REFERÊNCIAS
% =============================================================================
\newpage
\section*{Referências}
\addcontentsline{toc}{section}{Referências}

% As referências devem ser inseridas em um arquivo .bib no formato
% padrão BibTeX. Exemplo de entrada:
%
% @article{autor2023,
%   author  = {SOBRENOME, Nome and SOBRENOME2, Nome2},
%   title   = {Título do artigo},
%   journal = {Nome do Periódico},
%   year    = {2023},
%   volume  = {N},
%   number  = {N},
%   pages   = {XXX--XXX},
%   doi     = {10.XXXX/xxxxx}
% }
%
% @book{autor2020,
%   author    = {SOBRENOME, Nome},
%   title     = {Título do Livro},
%   edition   = {N},
%   address   = {Cidade},
%   publisher = {Editora},
%   year      = {2020}
% }

% Descomente as linhas abaixo e aponte para seu arquivo .bib:
% \bibliography{referencias}
% \bibliographystyle{abntex2-alf}   % estilo autor-data ABNT

% --- Lista de referências manuais (placeholder) ---
\begin{thebibliography}{99}

\bibitem[autor2023]{autor2023}
SOBRENOME, Nome; SOBRENOME, Nome. Título do artigo.
\textit{Nome do Periódico}, v.~N, n.~N, p.~XXX--XXX, 2023.
DOI: \href{https://doi.org/10.XXXX/xxxxx}{10.XXXX/xxxxx}.

\bibitem[autor2020]{autor2020}
SOBRENOME, Nome. \textbf{Título do Livro}. N. ed. Cidade: Editora, 2020.

\bibitem[autor2021]{autor2021}
NOME, Sobrenome et al. Título do capítulo. In:
ORGANIZADOR, Nome (Org.). \textbf{Título da Coletânea}.
Cidade: Editora, 2021. p.~XX--YY.

\bibitem[autor2022]{autor2022}
SOBRENOME, Nome. \textbf{Título da Tese}. 2022. N f.
Tese (Doutorado em [Área]) — Universidade, Cidade, 2022.

\bibitem[autorb2022]{norma2022}
BRASIL. \textbf{Lei nº 14.133, de 1º de abril de 2021}.
Lei de Licitações e Contratos Administrativos. Diário Oficial da
União, Brasília, DF, 2021.

% Adicione mais referências conforme necessário

\end{thebibliography}

% =============================================================================
% ANEXOS (OPCIONAIS)
% =============================================================================
\newpage
\appendix
\section*{Anexos}
\addcontentsline{toc}{section}{Anexos}

\subsection*{Anexo A — Currículo Resumido do Coordenador}
[Inclua um resumo curricular (máx. 2 páginas) destacando:
produção científica nos últimos 5 anos, orientações concluídas,
projetos financiados anteriormente, prêmios e distinções.]

\vspace{0.5cm}
\subsection*{Anexo B — Cartas de Anuência}
[Cópias das cartas de anuência dos colaboradores e instituições
coparticipantes, quando exigidas pelo edital.]

\vspace{0.5cm}
\subsection*{Anexo C — Documentação Complementar}
[Outros documentos solicitados pelo edital: declaração de
infraestrutura, comprovante de submissão ao CEP, comprovante de
vínculo institucional, etc.]

% =============================================================================
% FIM DO DOCUMENTO
% =============================================================================
\end{document}
